% Hodina 65, rozpis řádek 65. Zdroje: učebnice
% ucebnice/02_Funkce_a_rovnice/Linearni_funkce_a_rovnice/Lineární funkce.tex,
% oddíl Způsoby, jak vytvořit graf; Realisticky.cz lokální 020109 a 020110,
% lekce o lineární funkci – metodika konstrukce. Úlohy a graf vlastní.
\documentclass[11pt,a4paper]{article}
% Spolecny styl pro ucitelske pripravy MAT (Work/MAT)
% Verze 1.1, 2026-10-23
\NeedsTeXFormat{LaTeX2e}
\ProvidesPackage{priprava}[2026/10/23 v1.1 Sablona ucitelskych priprav MAT]

% Zakladni typografie a cestina
\RequirePackage[T1]{fontenc}
\RequirePackage[utf8]{inputenc}
\RequirePackage[czech]{babel}
\RequirePackage{lmodern}
\RequirePackage{microtype}
\RequirePackage{amsmath,amssymb}
\RequirePackage{xcolor}
\RequirePackage{enumitem}
\RequirePackage[most]{tcolorbox}
\RequirePackage{tabularx,booktabs,array}
\RequirePackage{fancyhdr}
\RequirePackage{needspace}
\RequirePackage{geometry}
\RequirePackage[hidelinks]{hyperref}

\geometry{a4paper,top=17mm,bottom=18mm,left=23mm,right=23mm}
\definecolor{notegray}{gray}{0.39}
\setlength{\parindent}{0pt}
\setlength{\parskip}{2.5pt plus 1pt minus 0.5pt}
\setlist[itemize]{topsep=2pt,itemsep=1pt,parsep=0pt,leftmargin=1.45em}
\setlist[enumerate]{topsep=2pt,itemsep=2pt,parsep=0pt,leftmargin=1.65em}
\setcounter{secnumdepth}{0}
\renewcommand{\arraystretch}{1.12}

% Zahlavi: vlevo rocnik | globalni tema, vpravo cislo: nazev hodiny
% Zapati zustava zcela prazdne.
\pagestyle{fancy}
\fancyhf{}
\renewcommand{\headrulewidth}{0.35pt}
\renewcommand{\footrulewidth}{0pt}
\setlength{\headheight}{14pt}
\newcommand{\nastavPripravu}[4]{%
  \fancyhead[L]{\small #1\enspace\textbar\enspace #2}%
  \fancyhead[R]{\small #3: #4}%
}

% Minimalisticke cernobile boxy. Kazda definice zvlast.
\tcbset{commonlesson/.style={enhanced,breakable=false,colback=white,colframe=black,
  colbacktitle=white,coltitle=black,boxrule=0.45pt,
  arc=0.8mm,outer arc=0.8mm,left=3mm,right=3mm,top=1.6mm,bottom=1.6mm,
  toptitle=1mm,bottomtitle=0.7mm,fonttitle=\bfseries,
  attach boxed title to top left={xshift=2.5mm,yshift=-1.15mm},
  boxed title style={colback=white,colframe=white,boxrule=0pt,left=1mm,right=1mm}}}
\newtcolorbox{zadaniBox}[1][]{commonlesson,title={Zadání},#1}
\newtcolorbox{prikladBox}[1][]{commonlesson,title={Příklad},#1}
\newtcolorbox{definiceBox}[1][]{commonlesson,title={Definice},#1}

% Spolecny box pro zadani a jeho reseni, bez mezery mezi castmi.
% Pouziti:
%   \begin{prikladReseniBox}
%     Zadani ulohy.
%     \tcblower
%     \begin{align*} ... \end{align*}
%   \end{prikladReseniBox}
% Spodni cast ma jemny sedy podklad a delici linku; lze v ni sazet
% vice radku matematiky. Obe casti maji jen jeden spolecny ramecek.
\tcbset{solutionlesson/.style={
  commonlesson,
  enhanced jigsaw,
  breakable,
  colbacklower=black!3,
  segmentation style={solid,draw=black!25,line width=0.35pt},
  fontlower=\small\color{notegray},
  before lower={\noindent\textbf{Řešení:}\enspace},
  leftlower=3mm,rightlower=3mm
}}
\newtcolorbox{prikladReseniBox}[1][]{solutionlesson,title={Příklad},#1}
\newtcolorbox{zadaniReseniBox}[1][]{solutionlesson,title={Zadání},#1}

% Pomucky pro autora pripravy: sedive poznamky; \reseni pro starsi pripravy
\newcommand{\poznamka}[1]{\par\begingroup\small\color{notegray}\itshape
  \textnormal{\textbf{Poznámka pro učitele.}}\enspace #1\par\endgroup}
\newcommand{\reseni}[1]{\par\begingroup\small\color{notegray}
  \textbf{Očekávané řešení:}\enspace #1\par\endgroup}
\newcommand{\blok}[1]{\needspace{5\baselineskip}\subsection*{#1}}
\newcommand{\vyrok}[1]{\uv{#1}}

\usepackage{tikz}
\nastavPripravu{1. ročník}{Funkce}{65}{Sestrojení grafu lineární funkce}
\begin{document}
\blok{Návaznost a princip}
Z hodiny 64 víme, že $b=f(0)$ určuje bod na svislé ose a směrnice $a=\Delta y/\Delta x$ určuje sklon. Dnes obojí využijeme při přesné konstrukci. U předpisu $f(x)=ax+b$ na $\mathbb R$ stačí dva různé body, protože přímku určují jednoznačně. Třetí vypočtený bod poskytne užitečnou kontrolu. Je-li funkce zadána jen na intervalu, kreslíme pouze příslušnou část přímky.
\blok{Konstrukce z tabulky}
Nejprve zvolíme vhodné argumenty, dosadíme do předpisu a vytvoříme tabulku. Na obou osách stanovíme stejnoměrné a popsané měřítko. Body $(x,f(x))$ zakreslíme a spojíme přímkou; délku grafu přizpůsobíme oboru.
\begin{prikladReseniBox}
Sestrojte graf $f(x)=\frac12x+1$ na $\mathbb R$ pomocí tabulky a ověřte bod pro $x=4$.
\tcblower
Volíme sudé argumenty, aby hodnoty vyšly celá čísla:
\[
\begin{array}{c|rrrr}x&-2&0&2&4\\\hline f(x)&0&1&2&3\end{array}
\]
Zakreslíme $(-2;0)$, $(0;1)$, $(2;2)$ a $(4;3)$. První dva body určují přímku; další dva kontrolují sazbu a výpočet. Dosazení $f(4)=2+1=3$ potvrzuje poslední bod.
\end{prikladReseniBox}
\begin{center}
\begin{tikzpicture}[x=.8cm,y=.8cm,>=stealth]
\draw[step=1,gray!25,very thin] (-4,-1) grid (5,4);
\draw[->] (-4.3,0)--(5.4,0) node[right] {$x$};\draw[->] (0,-1.3)--(0,4.3) node[above] {$y$};
\foreach \x in {-4,-2,2,4}{\draw (\x,.07)--(\x,-.07) node[below] {\small $\x$};}
\foreach \y in {1,2,3,4}{\draw (.07,\y)--(-.07,\y) node[left] {\small $\y$};}
\draw[very thick] (-4,-1)--(5,3.5);
\foreach \x/\y in {-2/0,0/1,2/2,4/3}{\fill (\x,\y) circle (2.2pt);}
\node[above left] at (-2,0) {\small $(-2;0)$};\node[below right=5pt] at (4,3) {\small $(4;3)$};
\end{tikzpicture}
\end{center}
\blok{Konstrukce pomocí směrnice a úseku}
Zaznačíme $P=(0;b)$. Zvolíme vodorovný krok $\Delta x\ne0$ a vypočteme svislý krok $\Delta y=a\Delta x$. Druhý bod je $Q=(\Delta x;b+a\Delta x)$. Při zlomkové směrnici bývá výhodné zvolit jmenovatel jako vodorovný krok.
\begin{prikladReseniBox}
Sestrojte $g(x)=-\frac23x+2$ pomocí bodu na ose $y$ a směrnice. Ověřte druhý bod dosazením.
\tcblower
První bod je $P=(0;2)$. Vezmeme $\Delta x=3$; pak $\Delta y=(-2/3)\cdot3=-2$. Dostaneme $Q=(3;0)$. Kontrola: $g(3)=-(2/3)\cdot3+2=0$. Grafem je přímka procházející $P,Q$; při pohybu doprava klesá.
\end{prikladReseniBox}
\begin{center}
\begin{tikzpicture}[x=.9cm,y=.75cm,>=stealth]
\draw[step=1,gray!25,very thin] (-1,-1) grid (4,3);
\draw[->] (-1.3,0)--(4.4,0) node[right] {$x$};\draw[->] (0,-1.3)--(0,3.4) node[above] {$y$};
\draw[thick] (-1,{8/3})--(4,{-2/3});
\draw[dashed] (0,2)--(3,2)--(3,0);
\fill (0,2) circle (2.3pt) node[below left=6pt] {\small $P(0;2)$};
\fill (3,0) circle (2.3pt) node[below left=5pt] {\small $Q(3;0)$};
\node[above] at (1.5,2) {\small $\Delta x=3$};
\node[right] at (3,1) {\small $\Delta y=-2$};
\end{tikzpicture}
\end{center}
\begin{prikladReseniBox}
Funkce $h(x)=2x-1$ je zadána jen pro $x\in\langle-1;2)$. Kterou část přímky nakreslíme a jak označíme konce?
\tcblower
Krajní hodnoty předpisu jsou $h(-1)=-3$ a $h(2)=3$. Kreslíme úsečku od plného bodu $(-1;-3)$ k prázdnému bodu $(2;3)$; pravý konec není součástí grafu. Kontrolní vnitřní bod $(0;-1)$ leží na úsečce.
\end{prikladReseniBox}
\begin{center}
\begin{tikzpicture}[x=.9cm,y=.55cm,>=stealth]
\draw[step=1,gray!25,very thin] (-2,-4) grid (3,4);
\draw[->] (-2.3,0)--(3.4,0) node[right] {$x$};\draw[->] (0,-4.3)--(0,4.4) node[above] {$y$};
\draw[very thick] (-1,-3)--(2,3);
\fill (-1,-3) circle (2.5pt) node[below left] {\small $(-1;-3)$};
\draw[fill=white,thick] (2,3) circle (2.8pt) node[above right] {\small $(2;3)$};
\fill (0,-1) circle (2pt);
\end{tikzpicture}
\end{center}
\poznamka{Kontroluj měřítko, pořadí souřadnic a vyznačení otevřeného konce. Při tabulce nepropojuj body přes mezeru v definičním oboru. Naopak u funkce na celé $\mathbb R$ nekončí graf v posledním tabulkovém bodě.}
\blok{Samostatné procvičení}
\begin{zadaniReseniBox}
\begin{enumerate}[label=\alph*)]
\item Pro $p(x)=2x-3$ vypočtěte body při $x=0,1,2$ a z nich načrtněte graf.
\item U $q(x)=\frac34x-2$ určete bod na ose $y$ a vhodný druhý bod pomocí kroku $\Delta x=4$.
\item U $r(x)=-x+1$ na $\langle0;3\rangle$ určete krajní body grafu.
\end{enumerate}
\tcblower
\begin{enumerate}[label=\alph*)]
\item Body $(0;-3),(1;-1),(2;1)$ leží na jedné přímce; rozdíl hodnot při kroku $1$ je vždy $2$.
\item $P=(0;-2)$, $\Delta y=(3/4)\cdot4=3$, tedy $Q=(4;1)$. Dosazení $q(4)=1$.
\item $r(0)=1$, $r(3)=-2$, tedy plné krajní body $(0;1)$ a $(3;-2)$.
\end{enumerate}
\end{zadaniReseniBox}
\blok{Volitelná rezerva}
\begin{zadaniReseniBox}
Dva žáci sestrojili graf $y=-\frac12x+3$. První použil body $(0;3),(2;2)$, druhý $(0;3),(-2;4)$. Mohou mít oba pravdu?
\tcblower
Ano. Dosazením $x=2$ dostaneme $2$, pro $x=-2$ dostaneme $4$. Všechny tři body leží na stejné přímce; vodorovný krok může být kladný i záporný, jen musí odpovídat svislému kroku podle směrnice.
\end{zadaniReseniBox}
\blok{Domácí práce}
\begin{zadaniReseniBox}
\begin{enumerate}[label=\alph*)]
\item Tabulkou pro $x=-1,0,1$ připravte graf $f(x)=-2x+1$.
\item Směrnicí sestrojte $g(x)=\frac25x-1$; použijte $\Delta x=5$.
\item U $h(x)=x+2$ na $(-1;2\rangle$ určete oba krajní body a jejich vyznačení.
\end{enumerate}
\tcblower
\begin{enumerate}[label=\alph*)]
\item Hodnoty jsou $3,1,-1$, tedy body $(-1;3),(0;1),(1;-1)$; leží na klesající přímce.
\item $P=(0;-1)$, $\Delta y=(2/5)\cdot5=2$, tedy $Q=(5;1)$; kontrola $g(5)=1$.
\item Při $x=-1$ by hodnota byla $1$, ale tento bod je prázdný. Při $x=2$ je hodnota $4$, bod $(2;4)$ je plný. Kreslíme úsečku mezi nimi.
\end{enumerate}
\end{zadaniReseniBox}
\poznamka{Jádro 45 min: připomenutí 3, konstrukce tabulkou 12, konstrukce směrnicí 10, omezený obor 5, samostatné kreslení 9, kontrola a organizace 6 min. Rezerva 5 min, domácí 10–15 min. Příští hodina probere vodorovnou přímku při nulové směrnici.}
\end{document}
